\begin{abstract}
LLM agents increasingly carry context forward through transcripts, summaries, and memory, making safety depend not only on what an agent outputs but on what the system stores and later reuses---including unsafe content from users, other agents, or prior history. Focusing on toxicity as a tractable axis of behavioral safety, we identify \textit{memory laundering}: a failure mode in which summarization removes explicit toxic language while preserving the conflict framing, hostile stance, or adversarial context that shapes later generations. Memory can thus appear safe to standard toxicity detectors while still steering downstream agents toward hostile responses. Using paired counterfactual multi-agent rollouts, we show that 100\% of toxic-condition memory summaries across 200 paired seeds fall below commonly used toxicity thresholds, yet conditioning on them substantially increases downstream toxicity relative to matched neutral baselines (0.168 in toxicity units, p < 0.001). To capture this hidden influence, we introduce the sub-threshold propagation gap (SPG): the downstream toxicity difference between toxic- and neutral-origin conditions, restricted to memory states a deployed monitor would label safe. SPG reveals two propagation routes: transcript reuse drives overt downstream toxicity, while compressed memory carries hidden sub-threshold influence. Sanitizing toxic state before memory summarization nearly eliminates this hidden gap, whereas cleaning only the completed summary leaves it largely intact. These findings yield a concrete design rule---sanitize before summarizing---and reframe agent safety as a state-control problem over evolving context. 
%future agents consume.
\end{abstract}
